% Zumdahl Chapter 5 free-response set -- 2 long (10) + 3 short (4) = 32 pts.
\documentclass{apchem-exam}

\apcunitword{Chapter}
\apcunit{5}
\apcunittitle{Gases}
\apcdoctitle{Free-Response Set}
\apcsubtitle{Zumdahl \S5.1--\S5.9 \textbullet\ 5 questions \textbullet\ 32 points \textbullet\ 50 minutes}

\begin{document}
\apctitleblock

\begin{apcnote}[Scope]
Covers \textbf{CED 3.4} (ideal gas law, \S5.2--5.5), \textbf{3.5}
(kinetic molecular theory, \S5.6--5.7) and \textbf{3.6} (deviation from
ideal behaviour, \S5.8--5.9). \S5.10 (atmospheric chemistry) is context,
not assessed content, and does not appear.

$R = \SI{0.08206}{\litre\atm\per\mole\per\kelvin}$. The vapour pressure of
water at \SI{25}{\celsius} is \SI{23.8}{\torr}.
\end{apcnote}

\examsection{II}{Free Response}{5 questions \textbullet\ 32 points \textbullet\ 50 minutes}{\calculatorok}

% =====================================================================
\frq{long}{10}{Zumdahl \S5.3--\S5.5 \textbullet\ CED 3.4 \textbullet\ gas stoichiometry and partial pressure}

Calcium carbonate decomposes on strong heating:
\[ \ce{CaCO3(s) -> CaO(s) + CO2(g)} \]
A \SI{10.0}{\gram} sample is decomposed completely and the carbon dioxide
is collected in a rigid \SI{2.00}{\litre} vessel at \SI{25}{\celsius}.

\begin{parts}
  \item Calculate the number of moles of \ce{CO2} produced.
        \answerlines[2]{$n(\ce{CaCO3}) = 10.0/100.09 = \num{0.0999}$~mol;
        the ratio is $1:1$, so $n(\ce{CO2}) = \mathbf{0.0999}$~mol}
  \item Calculate the pressure exerted by the \ce{CO2} in the vessel.
        \workspace[3cm]
        \keyonly{\begin{apcanswer}$T = 25 + 273 = \SI{298}{\kelvin}$.

        $P = \dfrac{nRT}{V} = \dfrac{(0.0999)(0.08206)(298)}{2.00}
        = \mathbf{1.22}$~atm\end{apcanswer}}
  \item The gas is instead collected over water at \SI{25}{\celsius},
        and the total pressure in the collection tube is
        \SI{755}{\torr}. Calculate the partial pressure of the \ce{CO2}.
        \answerlines[2]{$P(\ce{CO2}) = 755 - 23.8 = \mathbf{731}$~torr
        --- the collected gas is saturated with water vapour, which
        contributes its own partial pressure.}
  \item State what happens to the pressure in the rigid vessel if the
        temperature is raised to \SI{100}{\celsius}, and justify your
        answer using kinetic molecular theory. \workspace[2.8cm]
        \keyonly{\begin{apcanswer}The pressure \textbf{rises}. Volume and
        amount of gas are fixed, so by $P = nRT/V$ the pressure is
        proportional to absolute temperature: $373/298$ gives about
        \SI{1.53}{\atm}.

        In kinetic terms, raising the temperature raises the average
        kinetic energy and therefore the average speed of the molecules.
        They strike the walls both \emph{more often} and \emph{harder},
        and pressure is the force of those collisions per unit
        area.\end{apcanswer}}
  \item Explain why the identity of the gas does not appear anywhere in
        the ideal gas law. \answerlines[2]{The ideal gas model treats
        molecules as point particles with no volume and no intermolecular
        forces, so nothing in it depends on what the molecules are ---
        only on how many there are and how fast they move. Equal numbers
        of moles of any ideal gas exert the same pressure under the same
        conditions.}
\end{parts}

\begin{rubric}
  \rpoint{2}{(a) 1 pt moles of \ce{CaCO3}; 1 pt the $1:1$ ratio}
  \rpoint{2}{(b) 1 pt temperature converted to kelvin; 1 pt
    \SI{1.22}{\atm}. Using \SI{25}{} in place of \SI{298}{\kelvin} earns 0
    for the whole part}
  \rpoint{2}{(c) 1 pt subtracting the water vapour pressure; 1 pt
    \SI{731}{\torr}}
  \rpoint{2}{(d) 1 pt pressure rises, tied to $P \propto T$ at fixed $n$
    and $V$; 1 pt collisions become more frequent \emph{and} more
    forceful. ``The molecules move faster'' alone earns 1}
  \rpoint{2}{(e) 1 pt no intermolecular forces / negligible molecular
    volume in the model; 1 pt so behaviour depends on amount, not identity}
\end{rubric}

% =====================================================================
\frq{long}{10}{Zumdahl \S5.6--\S5.9 \textbullet\ CED 3.5, 3.6 \textbullet\ molecular speeds and real gases}

A sealed flask contains a mixture of helium (\SI{4.003}{\gram\per\mole})
and sulfur dioxide (\SI{64.07}{\gram\per\mole}) at the same temperature.

\begin{parts}
  \item State how the average kinetic energies of the two gases compare,
        and justify your answer. \answerlines[2]{They are \textbf{equal}.
        Average kinetic energy depends only on absolute temperature, and
        both gases are at the same temperature.}
  \item State how their average molecular speeds compare, and explain the
        difference from (a). \answerlines[3]{Helium is much faster. Since
        $\text{KE} = \tfrac{1}{2}mv^2$ is the same for both, the lighter
        molecule must have the greater speed. Equal energy does not mean
        equal speed --- mass is the other factor.}
  \item Calculate the ratio of the effusion rate of helium to that of
        sulfur dioxide. \workspace[2.6cm]
        \keyonly{\begin{apcanswer}$\dfrac{\text{rate}_{\ce{He}}}
        {\text{rate}_{\ce{SO2}}} = \sqrt{\dfrac{M_{\ce{SO2}}}
        {M_{\ce{He}}}} = \sqrt{\dfrac{64.07}{4.003}} = \sqrt{16.01}
        = \mathbf{4.00}$

        Helium effuses four times as fast.\end{apcanswer}}
  \item Sketch, in words, how the distribution of molecular speeds for
        helium differs in shape from that for sulfur dioxide at the same
        temperature. \answerlines[3]{The helium curve peaks at a much
        higher speed and is broader and lower --- its molecules are spread
        over a wide range of high speeds. The sulfur dioxide curve peaks
        at a low speed and is narrower and taller. Both enclose the same
        total area, because both represent the same number of molecules.}
  \item At high pressure the measured pressure of a real gas is
        \emph{lower} than the ideal gas law predicts. Explain why, in
        terms of intermolecular forces. \workspace[2.6cm]
        \keyonly{\begin{apcanswer}The ideal model assumes molecules do not
        attract one another. In a real gas they do. A molecule about to
        strike the wall is pulled back by attractions from the molecules
        behind it, so it hits with less force than it otherwise would.
        Every collision is slightly softened, and the measured pressure
        falls below the ideal prediction. The effect grows as pressure
        rises, because the molecules are closer together and the
        attractions matter more.\end{apcanswer}}
\end{parts}

\begin{rubric}
  \rpoint{2}{(a) 1 pt ``equal''; 1 pt because KE depends only on
    temperature}
  \rpoint{2}{(b) 1 pt helium faster; 1 pt an argument from
    $\text{KE} = \tfrac{1}{2}mv^2$ with equal KE. Saying only ``helium is
    lighter'' earns 1}
  \rpoint{2}{(c) 1 pt the inverse square-root relationship; 1 pt
    \num{4.00}. An inverted ratio (\num{0.25}) earns 1}
  \rpoint{2}{(d) 1 pt helium's peak at higher speed and broader; 1 pt
    equal areas / same number of molecules}
  \rpoint{2}{(e) 1 pt attractions between molecules; 1 pt reduced force of
    wall collisions. An answer resting on molecular volume describes the
    opposite deviation and earns 0}
\end{rubric}

% =====================================================================
\frq{short}{4}{Zumdahl \S5.2--\S5.3 \textbullet\ CED 3.4 \textbullet\ applying the gas laws}

A sample of nitrogen occupies \SI{4.00}{\litre} at \SI{1.50}{\atm} and
\SI{300}{\kelvin}.

\begin{parts}
  \item Calculate the volume it would occupy at \SI{1.00}{\atm} and
        \SI{450}{\kelvin}. \workspace[2.6cm]
        \keyonly{\begin{apcanswer}$\dfrac{P_1V_1}{T_1} =
        \dfrac{P_2V_2}{T_2}$

        $V_2 = 4.00 \times \dfrac{1.50}{1.00} \times \dfrac{450}{300}
        = \mathbf{9.00}$~L\end{apcanswer}}
  \item State whether the number of moles of nitrogen has changed, and why.
        \answerlines[2]{No. The sample is the same gas throughout; only
        its pressure, volume and temperature changed. $n$ is constant,
        which is exactly what allows the combined gas law to be used.}
\end{parts}

\begin{rubric}
  \rpoint{3}{(a) 1 pt correct combined-law setup; 1 pt both ratios applied
    in the right sense; 1 pt \SI{9.00}{\litre}}
  \rpoint{1}{(b) unchanged, because it is the same sealed sample}
\end{rubric}

% =====================================================================
\frq{short}{4}{Zumdahl \S5.5 \textbullet\ CED 3.4 \textbullet\ partial pressure and mole fraction}

A rigid vessel contains \SI{0.300}{\mole} \ce{N2}, \SI{0.100}{\mole}
\ce{O2} and \SI{0.100}{\mole} \ce{Ar}. The total pressure is
\SI{2.50}{\atm}.

\begin{parts}
  \item Calculate the partial pressure of nitrogen. \workspace[2.4cm]
        \keyonly{\begin{apcanswer}Total moles $= 0.300 + 0.100 + 0.100
        = \num{0.500}$.

        Mole fraction of \ce{N2} $= 0.300/0.500 = \num{0.600}$.

        $P(\ce{N2}) = 0.600 \times 2.50 = \mathbf{1.50}$~atm\end{apcanswer}}
  \item Argon has the largest molar mass of the three. State whether it
        therefore exerts the largest partial pressure, and justify your
        answer. \answerlines[3]{No. Partial pressure depends on the
        \emph{number} of particles, not their mass --- it is set by mole
        fraction alone. Argon and oxygen are present in equal amounts, so
        they exert equal partial pressures (\SI{0.500}{\atm} each), and
        nitrogen exerts the largest because there is most of it.}
\end{parts}

\begin{rubric}
  \rpoint{2}{(a) 1 pt mole fraction \num{0.600}; 1 pt \SI{1.50}{\atm}}
  \rpoint{2}{(b) 1 pt ``no''; 1 pt partial pressure depends on mole
    fraction, not molar mass}
\end{rubric}

% =====================================================================
\frq{short}{4}{Zumdahl \S5.8--\S5.9 \textbullet\ CED 3.6 \textbullet\ conditions for ideal behaviour}

\begin{parts}
  \item State the two conditions under which a real gas behaves most
        nearly ideally. \answerlines[2]{\textbf{High temperature} and
        \textbf{low pressure}.}
  \item Explain, for each condition, why the gas approaches ideal
        behaviour. \workspace[3cm]
        \keyonly{\begin{apcanswer}\textbf{Low pressure.} The molecules are
        far apart, so the volume they themselves occupy is a negligible
        fraction of the container, and they are too far apart for
        attractions to matter.

        \textbf{High temperature.} The molecules have large kinetic
        energies compared with the strength of the intermolecular
        attractions, so a passing attraction barely changes their motion
        and they behave as though no forces acted at all.\end{apcanswer}}
\end{parts}

\begin{rubric}
  \rpoint{2}{(a) 1 pt each condition}
  \rpoint{2}{(b) 1 pt low pressure linked to negligible molecular volume
    and/or wide separation; 1 pt high temperature linked to kinetic energy
    overwhelming attractions}
\end{rubric}

\examtotal

\end{document}
